%% ===================================================================
\section{Related Work}
\label{sec:related}
%% ===================================================================

We position \AEGIS against four research threads.

\textbf{Adversarial attacks on graph structure.}
Nettack~\cite{zugner2018adversarial} and Mettack~\cite{zugner2019adversarial} demonstrated structural vulnerability via targeted and global poisoning using GCN~\cite{kipf2017semi} surrogates. Subsequent work explored RL-based~\cite{dai2018adversarial}, spectral~\cite{bojchevski2019adversarial}, and bilevel~\cite{xu2019topology} formulations; Wu et al.~\cite{wu2019adversarial} survey the landscape. Gosch et al.~\cite{gosch2023adversarial} identify pitfalls in adversarial training for GNNs. These methods use surrogate gradients or heuristic search and provide per-edge gradient information (e.g., Mettack's meta-gradients yield per-edge vulnerability scores), but optimize surrogate objectives rather than directly characterizing the sensitivity geometry. \AEGIS instead computes the closed-form first-order maximally sensitive perturbation direction via SVD of $S_c$, which is provably optimal for equilibrium shift within the tangent plane.

\textbf{Certified robustness and defense.}
Convex relaxations~\cite{bojchevski2019certifiable,zugner2019certifiable}, randomized smoothing~\cite{bojchevski2020efficient,scholten2022randomized}, and deterministic certificates~\cite{li2025agnncert,geisler2021robustness} provide robustness guarantees. Schuchardt et al.~\cite{schuchardt2023localized} extend smoothing to localized per-node certificates, partially addressing the per-node granularity gap. These methods provide \emph{certification} (which nodes are guaranteed robust) but do not identify which specific \emph{edges} drive vulnerability or provide the globally most sensitive perturbation direction. Empirical defenses~\cite{zhu2019robust,jin2020graph,zhang2020gnnguard} harden models without formal vulnerability maps. \AEGIS complements these approaches by providing per-edge rankings, per-node first-order sensitivity thresholds, and the SVD-identified maximally sensitive direction---\emph{structural vulnerability differentiation} rather than robustness certification (\cref{sec:certificates}). The two are complementary: certification answers ``is this node safe?'' while \AEGIS answers ``which edges are the weakest, and in what direction?''

\textbf{Implicit networks and equilibrium analysis.}
The DEQ framework~\cite{bai2019deep} introduced fixed-point models; IGNN~\cite{gu2020implicit} enforces unique equilibria via spectral normalization~\cite{miyato2018spectral}; EIGNN~\cite{liu2021eignn} improves efficiency. For robustness, monotone operator networks~\cite{winston2020monotone} provide built-in Lipschitz bounds, and El Ghaoui et al.~\cite{el2021implicit} analyze well-posedness. All address input sensitivity $\norm{\partial \zstar / \partial x}$, not structural sensitivity $\norm{\partial Z / \partial A}$. \AEGIS addresses structural sensitivity via the resolvent $(I-J_z)^{-1} J_A$ for implicit models (\cref{thm:phase_transition}) and the unrolled Jacobian $S_K$ for explicit ones (\cref{prop:explicit}).

\textbf{Sensitivity analysis, Lipschitz bounds, and explainability.}
Gama et al.~\cite{gama2020stability} derive stability bounds for graph filters and GNNs under graph perturbation from a signal-processing perspective, providing worst-case Lipschitz-type bounds on output change. Kenlay et al.~\cite{kenlay2021stability} extend this to spectral perturbation bounds that yield per-node stability guarantees. Monotone operator equilibrium networks~\cite{winston2020monotone,revay2020lipschitz} provide built-in Lipschitz certificates that implicitly bound per-input sensitivity. These approaches provide \emph{worst-case scalar bounds} on output sensitivity, but do not yield the \emph{directional} information that \AEGIS provides: the SVD of $S_c$ identifies not just how much the output can change, but in which perturbation direction, enabling both vulnerability ranking and directed attack analysis. Recent work on graph curvature~\cite{topping2022understanding,digiovanni2023oversquashing} connects edge-level structural properties (Ricci curvature, over-squashing bottlenecks) to information flow; negatively curved edges are bottlenecks that may correlate with high $S_c$ vulnerability scores, though we do not test this connection empirically.
\AEGIS draws on matrix perturbation theory~\cite{stewart1990matrix,trefethen2005spectra}. GNN explanation methods (GNNExplainer~\cite{ying2019gnnexplainer}, PGExplainer~\cite{luo2020parameterized}, gradient-based~\cite{pope2019explainability}) optimize for \emph{prediction fidelity}, not adversarial vulnerability: they identify edges important for a correct prediction, while \AEGIS identifies edges whose perturbation maximally damages predictions. Our case study (\cref{sec:power_flow}) connects to GNN-based power flow~\cite{donon2019graph,bolz2019power,kim2025physicsinformed} and contingency screening~\cite{nakiganda2023graph}.

\textbf{Over-smoothing and depth.} Deep GNNs suffer from over-smoothing~\cite{chen2020simple}, which GRAND/GRAND++~\cite{chamberlain2021grand,thorpe2022grand} address via continuous-depth diffusion. Entezari et al.~\cite{entezari2020all} and Mujkanovic et al.~\cite{mujkanovic2022defenses} show that many adversarial attacks fail against simple defenses or are artifacts of evaluation. \AEGIS sidesteps this debate by analyzing the model's own sensitivity surface rather than searching for attacks via surrogates.
